\documentclass[conference]{IEEEtran}
\usepackage[utf8]{inputenc}
\usepackage[T1]{fontenc}
\usepackage{cite}
\usepackage{amsmath,amssymb}
\usepackage{graphicx}
\usepackage{booktabs}
\usepackage{url}

\title{Measuring Validator and Sandboxed-Execution Coverage for LLM‑Generated NetOps Artifacts: A Paired Confirmatory Study}

\author{
\IEEEauthorblockN{Autonomous AI Researcher}
\IEEEauthorblockA{CNSM 2026 Agentic AI Researcher Track}
}

\begin{document}

\maketitle

\begin{abstract}
We preregistered and executed a paired confirmatory experiment (study\_id f2c3d8b1-7a6e-4a9a-b2c3-5f8e1d2b6c7e) that measured the marginal detection benefit of adding sandboxed emulation to a deterministic validator for LLM-generated intent\ensuremath{\rightarrow}configuration artifacts. Under shared\_initial\_candidate semantics using the declared model "gpt-5-mini" (provider: openai\_responses) we evaluated Npairs = 720 confirmatory tasks. The deterministic validator (deterministic\_netops\_validator\_v5) flagged 719/720 = 0.9986111111111111; the guarded pipeline (validator + sandboxed emulation) flagged 720/720 = 1.0. Paired difference (guarded - baseline) = 0.001388888888888884; exact McNemar two-sided p = 1.0; 95\% bootstrap percentile CI (10000 resamples, seed = 1729) = [0.0, 0.004166666666666667]. Run accounting: model\_calls\_used = 721 (720 shared initial + 1 repair-stage call). These artifact-grounded results lie far below our preregistered minimum effect-of-interest (0.20); we report the preregistered design, exact quantitative outcomes, run-accounting diagnostics, artifact-grounded interpretation for the negligible guarded benefit in this regime, and reproducibility pointers into the archived execution bundle.
\end{abstract}

\section{Introduction}
Generative large language models (LLMs) are increasingly proposed to translate operator intent into device configuration and NetOps workflows. Operational deployment requires generated artifacts to satisfy syntactic correctness, workflow ordering, and higher-order safety/policy constraints. A common mitigation architecture is generate\ensuremath{\rightarrow}validate\ensuremath{\rightarrow}(repair)\ensuremath{\rightarrow}execute, sometimes augmented by sandboxed emulation to reveal runtime-only failures that static validators miss. Prior surveys and system reports advocate such workflow-centred mitigations [1]--[3], but provide limited large-scale paired measures of the marginal value added by sandboxed execution. Motivated by this gap, we preregistered and executed a controlled paired study to measure the aggregate detection-rate difference contributed uniquely by an automated sandboxed-emulation stage when the same initial generated candidate is analyzed by both conditions.

Research question and contributions. Our preregistered research question was: for synthetic intent\ensuremath{\rightarrow}configuration tasks under shared\_initial\_candidate semantics, what is the paired detection-rate difference (guarded - baseline) between (A) a deterministic validator-only pipeline and (B) the same deterministic validator followed by automated sandboxed emulation? Contributions: (1) a machine-readable preregistration and faithful autonomous execution; (2) an artifact-grounded paired analysis on Npairs = 720 with complete accounting; (3) an artifact-grounded interpretation explaining a near-zero marginal effect; (4) concise reproducibility pointers into the archived bundle and prescriptive boundary conditions where guarded stages are likely to matter.

\section{Related Work}
Workflow-centred mitigations and verification tools. Recent survey and architectural reports document the practical trend of combining LLM generation with deterministic validators, emulators, and repair loops to raise operational assurance for NetOps artifacts [1]. These works emphasize that correctness in intent\ensuremath{\rightarrow}configuration pipelines is multidimensional: syntactic parsing, command-ordering/workflow constraints, invariants (e.g., keep‑one‑uplink‑up), and dynamic runtime properties may each require different V\&V primitives. Tool-focused case studies highlight the evolution and practical limits of static verifiers: mature configuration-analysis systems have progressively expanded rule-sets and topology-aware checks but still exhibit gaps for runtime-dependent behaviors and stateful sequences [2]. That evolution shows why researchers and operators often layer an emulator or sandbox after static validation: emulation can exercise temporal behaviors, device state transitions, and interactions that are difficult to capture exhaustively as static rules.

Empirical and prototype evidence. Prototype systems that integrate generation with verification and execution-harnesses report heterogeneous gains: in some task regimes a verifier alone catches most dangerous errors, while in others emulation uncovers subtle runtime failures or environment-specific side effects that static rules miss [1], [3]. The practical implication drawn across this literature is that the marginal utility of a guarded (validator+emulation) pipeline depends on (a) validator coverage relative to the task distribution, (b) the fidelity of the emulator, and (c) the distributional properties of the input intents (e.g., constrained, parser-validated intents versus noisy natural-language requests). These prior empirical and system-level observations motivated a paired experimental design: by reusing the identical initial generated candidate across conditions (shared\_initial\_candidate semantics) we isolate only downstream guarded-stage effects (emulation observations or permitted repair invocations) rather than conflating them with generation variability.

Gaps and evaluation needs. While surveys and system reports argue for generate\ensuremath{\rightarrow}validate\ensuremath{\rightarrow}repair and for emulator-assisted validation as best practice for high-assurance NetOps, they do not provide large-scale, preregistered paired estimates of the marginal detection benefit under controlled generation semantics. The literature also documents important open questions that our study was designed to probe: how much additional coverage does sandboxed emulation provide when a strong deterministic validator already targets the same workflow patterns, and which task regimes (adversarial inputs, cross‑vendor CLI heterogeneity, or stateful multi-step procedures) are most likely to reveal nontrivial guarded-stage benefits? By situating our paired confirmatory design amid these prior findings, we provide an artifact-grounded quantitative datapoint addressing one narrow slice of the broader evaluation agenda advocated in [1]--[3].

\section{Methodology}
Preregistration and experimental design. We preregistered a confirmatory paired-binary design (study\_id f2c3d8b1-7a6e-4a9a-b2c3-5f8e1d2b6c7e). Primary estimand: paired\_success\_rate\_difference\_guarded\_minus\_baseline aggregated across confirmatory samples. Planned confirmatory sample: Npairs = 720 (planned episodes = 1440). Minimum effect-of-interest for power planning = 0.20 (20 percentage points). The preregistration specified inclusion/exclusion rules, missingness handling, multiplicity handling for exploratory secondary tests, and deterministic contamination controls.

Task bank, vendors, and inclusion. Confirmatory items were synthetic intent\ensuremath{\rightarrow}configuration tasks from intent\_configuration\_repair\_v1 covering three abstracted vendor-CLI families and three topology templates (leaf-spine, 3-node service chain, triangle). Inclusion required vendor-specific parser acceptance; the preregistered missingness policy permitted up to two deterministic reseeds for parser failures; irrecoverable parser failures would be deterministically excluded and replaced. In execution there were zero irrecoverable parser exclusions: complete\_pair\_count = 720 (see analysis/missingness\_summary.csv).

Model configuration and sampling semantics. Declared/requested model: "gpt-5-mini" (provider field: openai\_responses). The archived model configuration records max\_output\_tokens = 2000 and temperature = null (execution/model\_configuration.json, sha256 a40e96e212ab87da...). Important clarification (preregistered and enforced): temperature = null indicates the client did not set an explicit temperature parameter; null does not imply temperature=0 or provider-side determinism. We enforced shared\_initial\_candidate semantics: for each pair we obtained exactly one sampled initial candidate (shared) which both baseline and guarded pipelines analyzed. Therefore any between-condition difference can only arise from (i) guarded-stage emulation observations not encoded as a validator rule, or (ii) the allowed downstream repair-stage invocation(s) recorded in run accounting.

Validator and guarded pipeline implementation. Baseline: deterministic\_netops\_validator\_v5 performing full intent-constraint validation (syntax, command ordering/workflow, required-field presence, and policy-level constraints) and returning binary detected/not-detected per the preregistered scoring rule full\_intent\_constraint\_validation. Guarded: baseline validator followed by an automated sandboxed-emulation harness executing the candidate against an emulated instance of the declared topology and producing runtime observations. The repair policy permitted at most one repair-stage model call per task; run accounting records exactly one repair-stage call in the confirmatory execution.

Execution controls, provenance, and deterministic reconciliation. The adapter family used: hosted\_netops\_gvr\_v1. Execution manifest and deterministic reconciliation artifacts record completed\_episode\_count = 1440, complete\_pair\_count = 720, failed\_episode\_count = 0, and provider-call audit information including model\_calls\_used = 721 and stage\_counts: shared\_initial\_generation = 720; repair = 1. Provenance and contamination controls (deterministic hashing, artifact-version checks) were applied per the preregistration; no contamination flags occurred for the confirmatory set.

\section{Results}
Primary confirmatory counts and test statistics (artifact-grounded). Analysis artifacts (analysis/paired\_contingency\_table.csv and analysis/condition\_summary.csv) report the contingency counts: n11 = 719 (detected by both), n10 = 1 (guarded-only), n01 = 0 (baseline-only), n00 = 0 (neither). Marginal detection rates: baseline = 719/720 = 0.9986111111111111; guarded = 720/720 = 1.0. Paired estimate (guarded - baseline) = 0.001388888888888884. Exact McNemar two-sided test p = 1.0. Paired-bootstrap percentile 95\% CI (10000 resamples, bootstrap\_seed = 1729) = [0.0, 0.004166666666666667]. All numeric values below are reproduced verbatim from analysis/results.json and analysis artifacts.

Table 1 --- Primary confirmatory summary (artifact-grounded).
- Npairs = 720 (complete paired episodes)
- Baseline successes = 719; Baseline success rate = 0.9986111111111111
- Guarded successes = 720; Guarded success rate = 1.0
- Paired difference (guarded - baseline) = 0.001388888888888884
- Exact McNemar two-sided p = 1.0
- 95\% bootstrap percentile CI = [0.0, 0.004166666666666667] (10000 resamples, seed = 1729)

Discordant-pair attribution and repair accounting. Provider-call accounting and deterministic reconciliation identify a single repair-stage invocation (stage\_counts: repair = 1) and a unique n10 discordant count. Because shared\_initial\_candidate semantics were enforced, the permissible causes for that single guarded-only detection are: (i) a guarded-stage emulation observation not encoded as a validator violation, or (ii) the single recorded repair-stage invocation. The confirmatory execution bundle contains the raw per-episode records; the unique discordant record can be located by searching execution/raw\_results.jsonl (input-results SHA-256 = df367ecc\allowbreak{}40a6f0ea\allowbreak{}f4021292\allowbreak{}60e11d84\allowbreak{}50884a4e\allowbreak{}a5fdb94a\allowbreak{}602e0d25\allowbreak{}13aef0e8) for the single entry where baseline\_score = 0 and guarded\_score = 1. Deterministic reconciliation is archived as analysis/deterministic\_reconciliation.json (sha256 49731a2e\allowbreak{}52e049cc\allowbreak{}f5b53c35\allowbreak{}ddac2351\allowbreak{}ad75e7e9\allowbreak{}5504dc4a\allowbreak{}42bcb096\allowbreak{}d0dec3aa) and its provider\_call\_audit confirms model\_calls\_used = 721 and repair = 1. The raw\_results record for the discordant entry contains the pair\_id and the validator\_trace and repair\_provider\_trace fields that link the guarded detection to the documented repair-stage call; these artifacts are part of the canonical execution bundle referenced by the execution manifest.

Representative per-episode diagnostics. Representative scoring traces archived under execution/scoring/task-000001-*.json ... task-000006-*.json include validator\_trace objects with normalized\_configuration, validation\_before/after, parsed\_command\_count, required\_command\_count, violation\_count, and validator\_version = 5.0. Those representative artifacts illustrate that the vast majority of shared\_initial candidates complied with validator rules (violation\_count = 0) and that sandboxed emulation largely corroborated validator findings in this curated task bank.

\section{Discussion}
Why was the guarded-stage aggregate benefit negligible? The archived evidence supports three artifact-grounded factors.

1) Validator ceiling effect. Deterministic\_netops\_validator\_v5 covered the syntactic and explicit semantic properties encoded in the synthetic task bank (shutdown--configure--restore, make-before-break uplink switchover, paired MTU changes). Baseline success rate \ensuremath{\approx} 0.9986 indicates near-universal validator coverage for these curated patterns, leaving negligible headroom for the guarded stage.

2) Task curation and inclusion policy. Confirmatory items were parser-validated synthetic tasks specifically designed to exercise workflow policies but not constructed adversarially to evade validator checks. This curation concentrates cases inside the validator's coverage envelope and reduces the prevalence of subtle runtime-only failure modes that emulation might reveal.

3) Shared\_initial\_candidate semantics and limited repair activity. Enforcing one sampled initial candidate per pair isolates guarded-stage contributions but removes between-condition generation variability. Provider\_call\_audit and deterministic\_reconciliation.json record one repair-stage invocation aligned with the single n10 discordant count; thus the observed discordance is attributable to downstream guarded-stage activity recorded in run accounting rather than generation differences.

Operational implications and boundary conditions. For environments where a high-coverage static validator exists and inputs match its expected patterns (parser-validated, constrained intents), adding automated sandboxed emulation may provide negligible aggregate additional detection at scale. This conditional conclusion must not be generalized: in operational regimes with noisier natural-language intents, adversarial inputs, broader vendor diversity, stateful multi-step interactions, or validator gaps, guarded runtime checks and higher-fidelity emulation are likely to provide substantive benefit. The archived per-episode traces support targeted follow-up hypotheses: adversarial task injection, broader vendor-CLI coverage, and complementary validator families should increase the guarded-stage marginal contribution.

\section{Limitations}
\begin{itemize}
\item Synthetic task bank and deterministic task generator produce limited ecological diversity and create a validator-ceiling effect; results may not generalize to noisier, adversarial, or production distributions.
\item Single declared/requested model family (gpt-5-mini) and single validator implementation (deterministic\_netops\_validator\_v5) limit generalizability across models and validator families.
\item Preregistered shared\_initial\_candidate semantics (one initial generation reused across paired conditions) isolate guarded-stage contribution but remove between-condition generation variability; this design choice constrains discordance sources to downstream guarded-stage observations or repair calls.
\item Sandboxed-emulation fidelity relative to vendor/OS production behaviors and large-scale stateful interactions is limited by the emulation harness used; higher-fidelity emulators could change outcomes in other regimes.
\item Study power planning targeted a minimum detectable paired difference of 0.20; the observed aggregate effect (\textasciitilde{}0.0014) lies far below that design sensitivity and should be interpreted as a narrowly scoped near-zero finding for this experimental regime rather than a general negative result for all NetOps workflows.
\item Provider-side nondeterminism and hidden caching remain residual uncertainties despite deterministic reconciliation procedures; these are documented in analysis/deterministic\_reconciliation.json and provider\_call\_audit artifacts.
\end{itemize}

\section{Conclusion}
We preregistered and executed a paired confirmatory study comparing a strong deterministic validator with the same validator augmented by sandboxed emulation on a curated synthetic intent\ensuremath{\rightarrow}configuration task bank. Over 720 paired tasks the guarded pipeline produced a negligible aggregate detection gain (0.001388888888888884) with exact McNemar p = 1.0 and 95\% bootstrap CI [0.0, 0.004166666666666667]. Run accounting shows one repair-stage invocation corresponding to the single guarded-only detection; because shared\_initial\_candidate semantics were enforced, archived per-episode traces permit independent verification of that mapping via execution/raw\_results.jsonl and analysis/deterministic\_reconciliation.json. We interpret the result as arising from validator ceiling effects and curated task design; follow-up work should focus on adversarial and production-like task banks, multi-validator ensembles, and higher-fidelity emulation to identify regimes where guarded stages add substantive operational value.

Reproducibility pointers (compact). The canonical execution bundle archived by the run contains: execution/execution\_manifest.json (execution summary and preregistration SHA-256), execution/raw\_results.jsonl (input-results SHA-256 df367ecc\allowbreak{}40a6f0ea\allowbreak{}f4021292\allowbreak{}60e11d84\allowbreak{}50884a4e\allowbreak{}a5fdb94a\allowbreak{}602e0d25\allowbreak{}13aef0e8), analysis/deterministic\_reconciliation.json (sha256 49731a2e\allowbreak{}52e049cc\allowbreak{}f5b53c35\allowbreak{}ddac2351\allowbreak{}ad75e7e9\allowbreak{}5504dc4a\allowbreak{}42bcb096\allowbreak{}d0dec3aa), execution/model\_configuration.json (sha256 a40e96e2\allowbreak{}12ab87da\allowbreak{}251ac0d6\allowbreak{}dbb75bc5\allowbreak{}43afa134\allowbreak{}38b5a6b9\allowbreak{}ebd07359\allowbreak{}7ffdd820), analysis/paired\_contingency\_table.csv (sha256 efe6e8e7\allowbreak{}016fa843\allowbreak{}a8226e91\allowbreak{}1dbde829\allowbreak{}e943903d\allowbreak{}719101da\allowbreak{}8d956b3f\allowbreak{}f899133c), and analysis/condition\_summary.csv (sha256 394bc690\allowbreak{}671b7311\allowbreak{}94f9b42b\allowbreak{}4dcbe28f\allowbreak{}ded6e2ec\allowbreak{}8cf898b7\allowbreak{}98312253\allowbreak{}43c22c28). The discordant episode is the unique raw\_results.jsonl entry with baseline\_score = 0 and guarded\_score = 1; that record contains the pair\_id and validator/emulation traces that map it to the single repair-stage invocation recorded in deterministic\_reconciliation.json. These artifacts are archived as the canonical reproducibility bundle referenced by the execution manifest and the preregistration record. (See Disclosure for exact manifest references.)

\section*{Disclosure Statement}
Disclosure: Autonomous post-lock research run executed under the frozen master prompt supplied to the system. Declared/requested model: "gpt-5-mini" (provider recorded as openai\_responses). Deterministic validator: deterministic\_netops\_validator\_v5. Orchestration adapter family: hosted\_netops\_gvr\_v1. Preregistration id: f2c3d8b1-7a6e-4a9a-b2c3-5f8e1d2b6c7e (preregistration SHA-256: 0169919f\allowbreak{}b8c5aedb\allowbreak{}2c5a78e0\allowbreak{}31f1ab5c\allowbreak{}3aeda405\allowbreak{}cde1ae80\allowbreak{}fb2f8647\allowbreak{}68595c1b; recorded in the execution manifest). Initial master-prompt immutable reference: provenance/master\_prompt.txt, SHA-256 = 1872df1e\allowbreak{}1805d2d9\allowbreak{}6940456c\allowbreak{}a016bd66\allowbreak{}5d1d5196\allowbreak{}add77f5a\allowbreak{}cdf1582b\allowbreak{}b39b15ba. Execution summary (artifact-grounded): completed\_episode\_count = 1440, complete\_pair\_count = 720, model\_calls\_used = 721 (720 shared initial + 1 repair-stage call); stage\_counts and provider-call audit are archived in analysis/deterministic\_reconciliation.json (sha256 49731a2e52e0...). Human role and autonomy boundary: a human Research Director selected the broad NetOps topic family and constrained the pre-lock master prompt; after the autonomy lock the agentic pipeline autonomously performed literature selection, preregistration, code and prompt generation, experiment execution, analysis, peer review, and manuscript drafting. No manual human editing of technical manuscript content occurred after the autonomy lock. The execution manifest and archived artifacts named above serve as the canonical reproducibility bundle.

\begin{thebibliography}{99}
\bibitem{ref1} Muhammad Bilal, Jon Crowcroft, Ruizhi Wang, Xiaolong Xu, Schahram Dustdar, "Large Language Models for Agentic NetOps and AIOps: Architectures, Evaluation, and Safety", 2026, 10.48550/arxiv.2605.12729
\bibitem{ref2} Matt Brown, Ari Fogel, Daniel Halperin, Victor Heorhiadi, Ratul Mahajan, Todd Millstein, "Lessons from the evolution of the Batfish configuration analysis tool", 2023, 10.1145/3603269.3604866
\bibitem{ref3} http://arxiv.org/abs/2407.08249
\end{thebibliography}

\end{document}
